\chapter*{서문}
\addcontentsline{toc}{chapter}{서문}

이 책은 VLA를 처음 배우는 일반 AI 독자를 위한 입문서가 아니다. 이 책은 로봇공학, 제어, 로봇러닝, 컴퓨터비전 배경을 가진 대학원생과 연구자가 Vision-Language-Action system을 engineering 관점에서 해석하기 위한 field manual이다.

이 책의 중심 질문은 ``어떤 모델이 가장 큰가''가 아니라 ``언어 조건부 perception과 learned action이 실제 로봇의 state estimation, controller, safety monitor, data engine, evaluation protocol과 어떤 책임 경계로 연결되는가''이다. 따라서 이 책은 model leaderboard를 만들지 않는다. 대신 VLA 논문과 시스템을 읽을 때 반드시 물어야 할 interface, evidence, failure mode를 정리한다.

\section*{이 책을 만든 이유}
\addcontentsline{toc}{section}{이 책을 만든 이유}

한 줄로 말하면, 이 책은 VLA를 ``모델 이름''이 아니라 ``로봇 시스템의 책임 구조''로 읽기 위해 만들었다.

최근 VLA 논문과 데모는 빠르게 늘고 있다. OpenVLA, RT 계열, $\pi_0$, FAST, humanoid foundation model, robot data engine, evaluation harness는 모두 중요한 흐름이다. 그러나 이 흐름을 backbone, parameter count, benchmark score 중심으로만 읽으면 로봇공학자가 실제로 물어야 할 질문이 뒤로 밀린다. 로봇은 문장을 맞히는 시스템이 아니라, 카메라와 센서로 불완전한 세계를 보고, 물체와 접촉하고, actuator limit 안에서 움직이며, 실패하면 멈추거나 복구해야 하는 물리 시스템이다.

따라서 VLA에서 진짜 어려운 질문은 다음과 같다.

\begin{quote}
VLM 또는 VLA policy가 만든 action은 어떤 state estimate, perception evidence, controller interface, safety constraint, data record, evaluation protocol을 통해 실제 robot behavior가 되는가?
\end{quote}

이 책은 이 질문을 앞에 둔다. 기존 VLA survey는 대개 model architecture, dataset, benchmark, representative system을 정리한다. 그런 정리는 필요하지만 충분하지 않다. 로봇공학자의 관점에서는 다음이 더 중요하다.

\begin{itemize}
  \item image가 곧 state가 아니라면, perception이 제공하던 depth, pose, mask, map, uncertainty 책임은 VLA 안에서 어디로 갔는가?
  \item policy-level action $\act_t$와 controller command $\ctrlcmd_t$는 왜 다르며, 그 경계가 논문 claim을 어떻게 바꾸는가?
  \item action token, delta pose, waypoint, action chunk, diffusion/flow action expert는 단순 output format인가, 아니면 control boundary인가?
  \item robot data는 web data와 무엇이 다르며, failure, correction, intervention, reset, controller metadata가 왜 필요한가?
  \item benchmark success는 capability의 충분한 증거인가, 아니면 특정 protocol 아래의 제한된 evidence인가?
  \item safety는 prompt policy인가, 아니면 perception, control, runtime monitor, data log, evaluation이 함께 만드는 stack property인가?
\end{itemize}

이 책의 차별점은 최신 VLA 이름을 많이 나열하는 데 있지 않다. 차별점은 VLA를 다음과 같은 system contract로 읽는 데 있다.

\begin{equation*}
  \text{Real VLA claim}
  =
  \{\text{model},\text{perception},\text{action},\text{control},\text{data},\text{evaluation},\text{safety}\}.
\end{equation*}

이 항목 중 하나가 비어 있으면 claim은 약해진다. 강한 VLA system은 model architecture만 제시하지 않는다. 어떤 observation을 받는지, 무엇을 state로 보거나 숨기는지, action이 어떤 단위와 frame과 rate를 갖는지, controller가 무엇을 보정하는지, failure가 data engine으로 어떻게 돌아가는지, evaluation이 무엇을 증명하고 무엇을 증명하지 못하는지, safety layer가 어떤 unsafe success를 막는지 설명해야 한다.

그래서 이 책은 classical robotics와 foundation model을 경쟁 관계로 놓지 않는다. 오히려 이 책의 출발점은 그 반대다. VLA가 강력해질수록 기존 robotics의 책임 일부는 network 내부로 흡수된다. 그러나 흡수된 책임은 사라지지 않는다. 그것은 logging, stress test, safety monitor, dataset card, evaluation report를 통해 다시 검증되어야 한다. 이 책은 바로 그 책임의 이동을 추적하는 책이다.

독자가 이 책에서 가져가야 할 것은 특정 모델의 순위가 아니다. 독자가 가져가야 할 것은 VLA 논문을 읽을 때 다음처럼 묻는 습관이다.

\begin{quote}
이 system은 무엇을 명시적 interface로 남겼고, 무엇을 hidden representation 안으로 넣었으며, 숨겨진 책임을 어떤 evidence로 다시 검증했는가?
\end{quote}

이 질문을 붙잡으면 VLA는 더 큰 VLM이 아니라 로봇 실행 stack의 문제로 보인다. 이것이 이 책을 만든 이유다.

\section*{이 책의 독자}
\addcontentsline{toc}{section}{이 책의 독자}

이 책은 다음 독자를 기준으로 한다.

\begin{itemize}
  \item 로봇공학, 제어, motion planning, robot learning 중 적어도 하나의 배경을 가진 대학원생
  \item VLM/LLM 기반 embodied AI 논문을 실제 robot deployment 관점에서 읽고 싶은 연구자
  \item VLA model을 사용하거나 fine-tuning할 때 action representation, controller interface, evaluation report를 설계해야 하는 엔지니어
\end{itemize}

이 책은 classical robotics textbook, reinforcement learning textbook, transformer textbook을 대체하지 않는다. 필요한 개념은 다시 정의하지만, 목적은 개념 자체의 교과서적 완결성이 아니라 Real VLA system claim을 검토하는 기준을 만드는 것이다.

\section*{이 책에서 다루지 않는 것}
\addcontentsline{toc}{section}{이 책에서 다루지 않는 것}

\begin{itemize}
  \item VLM/LLM 일반론 전체를 설명하지 않는다.
  \item 강화학습과 최적제어의 표준 교과서 내용을 처음부터 증명하지 않는다.
  \item 특정 VLA 모델의 leaderboard 순위를 고정하지 않는다.
  \item 실제 배치 evidence 없이 architecture novelty만으로 capability를 판정하지 않는다.
\end{itemize}

\section*{어떻게 읽을 것인가}
\addcontentsline{toc}{section}{어떻게 읽을 것인가}

Part I은 Real VLA가 올라갈 robot substrate를 고정한다. Part II는 imitation, RL, visuomotor policy, VLM grounding을 VLA-specific failure 관점에서 읽는다. Part III는 planner와 robotics transformer에서 action representation, data engine, adaptation으로 이어지는 현대 VLA의 중심부이다. Part IV는 evaluation, hierarchy, real-time, safety를 통해 deployment evidence를 다룬다. Part V는 humanoid, whole-book synthesis, future direction으로 범위를 넓힌다.

각 장을 읽을 때는 다음 네 질문을 유지한다.

\begin{enumerate}
  \item 이 장에서 바뀌는 interface는 무엇인가?
  \item 이 장의 주장은 어떤 변수, constraint, data record, metric으로 검증되는가?
  \item running example인 ``머그컵을 drying rack에 올리고 유리컵은 건드리지 말라''에서 어떤 failure를 설명하는가?
  \item 논문을 읽을 때 어떤 evidence가 없으면 claim을 약하게 봐야 하는가?
\end{enumerate}

\chapter*{Notation and Glossary}
\addcontentsline{toc}{chapter}{Notation and Glossary}

\section*{주요 notation}
\addcontentsline{toc}{section}{주요 notation}

\begin{center}
\small
\begin{tabularx}{\textwidth}{@{}p{0.18\textwidth}p{0.76\textwidth}@{}}
\toprule
기호 & 이 책에서의 의미 \\
\midrule
$\obs_t$ & time $t$의 observation. RGB, depth, tactile, proprioception, language context를 포함할 수 있다. \\
$\state_t$ & 실제 physical state. object pose, robot configuration, contact state, hidden disturbance 등을 포함한다. \\
$\bel_t$ & state uncertainty를 포함한 belief 또는 internal estimate. \\
$\task$ & language instruction 또는 task specification. \\
$\act_t$ & VLA policy가 내는 policy-level action. token, vector, waypoint, chunk, distribution일 수 있다. \\
$\ctrlcmd_t$ & controller, driver, actuator가 실행하는 control command. velocity, torque, impedance target 등이다. \\
$\conf_t$ & robot joint configuration 또는 generalized coordinate. \\
$C_t$ & constraint set. collision, contact, joint limit, safety envelope를 포함한다. \\
$h_t$ & history 또는 memory state. observation/action/outcome의 temporal context이다. \\
\bottomrule
\end{tabularx}
\end{center}

\section*{Notation discipline}
\addcontentsline{toc}{section}{Notation discipline}

이 책에서 반복되는 기호는 전역 의미를 우선한다. 특히 $\act_t$는 VLA policy의 policy-level action, $\ctrlcmd_t$는 controller command, $\conf_t$는 robot configuration으로 고정한다. symbolic planner의 discrete action은 필요하면 $\alpha_k$, skill은 $\sigma_i$, runtime record는 $\rho_t$, metadata는 $m_i$, contact mode는 $\mu_t$로 쓴다. 같은 문자 하나로 state, skill, query, reward, record를 동시에 나타내지 않는 것이 원칙이다.

일부 수식은 theorem이 아니라 diagnostic notation이다. 예를 들어 claim tuple, evidence vector, executability gap은 system을 읽기 위한 구조화된 표기이지, 모든 장에서 동일하게 측정되는 scalar metric이 아니다. 그런 경우 본문에서 ``diagnostic notation'' 또는 ``설계 원칙을 나타내는 notation''이라고 명시한다.

\section*{주요 용어}
\addcontentsline{toc}{section}{주요 용어}

\begin{center}
\small
\begin{tabularx}{\textwidth}{@{}p{0.24\textwidth}p{0.70\textwidth}@{}}
\toprule
Term & 이 책에서의 의미 \\
\midrule
VLM & image/text semantic model. 장면을 설명하거나 언어 질의를 처리하지만, 그 자체가 robot execution을 보장하지는 않는다. \\
VLA policy & observation, language, state estimate로부터 policy-level action을 내는 learned policy. \\
Real VLA system & VLA policy, controller, monitor, safety, data, evaluation까지 포함한 책임 구조. \\
Action representation & $\act_t$의 형식. token, continuous vector, waypoint, chunk, diffusion/flow sample 등이 포함된다. \\
Control boundary & $\act_t$가 $\ctrlcmd_t$로 바뀌는 경계. 이 경계가 불명확하면 논문 claim이 실제 robot behavior로 해석되지 않는다. \\
Data engine & data collection, curation, labeling, failure mining, evaluation replay가 반복되는 system loop. \\
Safety shield & action 또는 controller reference를 safe set으로 projection, reject, revalidate하는 runtime layer. \\
\bottomrule
\end{tabularx}
\end{center}
